\documentclass[10pt,a4paper]{amsart}
\usepackage[margin=19mm]{geometry}
\usepackage{amsmath,amssymb,mathtools}
\usepackage[scr=rsfs,cal=euler]{mathalfa}
\usepackage{xfrac}
\usepackage[unicode,hidelinks,bookmarks=false]{hyperref}
\newcommand{\ie}{i.e.\ }
\newcommand{\cf}{cf.\ }
\newcommand{\resp}{respectively\ }
\newcommand{\Subsection}{\subsection}
\providecommand{\nobreakdash}{\nobreak\hbox{-}}
\setlength{\emergencystretch}{2em}
\multlinegap0pt
\usepackage{bm}
\usepackage{bbm}
\usepackage{dsfont}
\usepackage{xspace}
\usepackage{cancel}
\usepackage{mathtools}
\usepackage{cleveref}
\usepackage{amsthm}
\makeatletter
\@ifundefined{thm}{\newtheorem{thm}{Thm}}{}
\@ifundefined{example}{\newtheorem{example}[thm]{Example}}{}
\@ifundefined{prop}{\newtheorem{prop}[thm]{Proposition}}{}
\makeatother
\DeclareMathOperator{\ASM}{{\tt ASM}}
\DeclareMathOperator{\Perm}{Perm}
\DeclareMathOperator{\codim}{{\tt codim}}
\DeclareMathOperator{\rk}{rk}
\title[Some algebraic properties of ASM varieties]{Some algebraic properties of ASM varieties}
\author{Ilani Axelrod-Freed, Hanson Hao, Matthew Kendall, Patricia Klein,, Yuyuan Luo}
\date{Source: arXiv:2505.10480v2 (CC BY 4.0)}
\begin{document}
\maketitle

\noindent\emph{Source and licence.} Some algebraic properties of ASM varieties. Version arXiv:2505.10480v2 (CC BY 4.0), distributed under the Creative Commons Attribution 4.0 International licence, as recorded in the arXiv licence metadata for this exact version and shown on the abstract page https://arxiv.org/abs/2505.10480v2. Full rights notice: licenses/arxiv-2505.10480-CC-BY-4.0.txt.\par\smallskip

\noindent\textit{Editorial scope.} Counted unit: the Prop whose source label is \texttt{prop:permBij} in the article (source0.tex lines 257--260, source label \texttt{prop:permBij}), delivered together with the article's own complete original proof of it (source0.tex lines 261--294, one \texttt{proof} environment of 34 lines). No mathematics is written, repaired or re-derived: statement and proof are the article's own text line for line. Required context retained uncounted: ex:permBij (lines 339--352). Every definition, display and label of those lines is retained unchanged. Omitted from this package and not redistributed: 1--256, 295--338, 353--1204. Nothing inside a retained excerpt is omitted or altered: every retained body is delivered exactly as the article's lines are written, so no omission receipt of any kind accompanies this package. Citations: the retained excerpts carry no citation command at all, so no bibliography entry of the article is transcribed and no reference is redistributed. Reference adaptation: none was needed and none was made; every reference inside the delivered unit resolves inside it, so no unresolved reference or citation is produced. Printed slips kept literal: no printed word, symbol, hypothesis or number of the counted statement or of its proof was altered. Layout and class: the article's own class is \texttt{amsart} with packages including inputenc, amsmath, amssymb, amsthm, verbatim, hyperref, xcolor, graphicx, caption, subcaption, mathtools, tikz, ytableau, soul; the wrapper is the lane's pinned \texttt{amsart} editorial wrapper at 10pt on A4, which restates the theorem-like environments the retained text needs and adds the article's own macro definitions that the retained text needs (\texttt{\textbackslash ASM}, \texttt{\textbackslash Perm}, \texttt{\textbackslash codim}, \texttt{\textbackslash rk}), with extra wrapper packages (bm, bbm, dsfont, xspace, cancel, mathtools, cleveref). The delivered numbering is the unit's own. Assets: no retained span contains a figure environment, a graphics-inclusion command, a PDF-page inclusion, a table or a TikZ picture; no image file is copied into this package, the package declares no asset dependency and no image file is redistributed with it. Licence: version arXiv:2505.10480v2 of this article is distributed under the Creative Commons Attribution 4.0 International licence, as recorded in the arXiv licence metadata for this exact version and shown on the abstract page https://arxiv.org/abs/2505.10480v2. A full rights notice is delivered at licenses/arxiv-2505.10480-CC-BY-4.0.txt. Characters: every retained excerpt is pure ASCII, so the delivered engine, XeLaTeX with its default Latin Modern text font, reports no missing character.
\begin{example}\label{ex:permBij}
 \Cref{prop:permBij} involves inserting into an ASM a row and column whose entries are $0$ except where they intersect, where the value is $1$. In the case of \Cref{prop:permBij}, it is specifically the first row and first column.  It is worth noting that equidimensionality is \emph{not} preserved by an arbitrary insertion of this type.  For example, consider $A = \begin{pmatrix}
0 & 1 & 0 \\
1 & -1 & 1 \\
0 & 1 &  0
\end{pmatrix}$ and $B = \begin{pmatrix}
0 & 1 & 0 & 0 \\
0 & 0 & 1 & 0 \\
1 & -1 & 0 & 1 \\
0 & 1 & 0 & 0
\end{pmatrix}$, in which case $B$ may be obtained from $A$ by inserting $(0,1,0,0)^T$ to become column $3$ and $(0,0,1,0)$ to become row $2$.

Then $I_A = (z_{1,1}, z_{1,2}z_{2,1}) = I_{312} \cap I_{231}$ note only defines an equidimensional ASM variety but even complete intersection while $I_B = (z_{1,1}, z_{2,1}, z_{1,2}z_{3,1}, z_{2,2}z_{3,1}) = I_{3412} \cap I_{2341}$ has one component of codimension $4$ and one of codimension $3$.  
\end{example}

\begin{prop}\label{prop:permBij}
Let $A \in \ASM(n)$ and $\Perm(A) =\{w_1,\ldots,w_k\}$.  For $w \in S_n$, the assignment $w \mapsto 1 \oplus w$ gives a bijection between $\Perm(A)$ and $\Perm(1 \oplus A)$.  In particular, $
\codim(A)=\codim(1 \oplus A)$, and $A$ is equidimensional if and only if $1 \oplus A$ is equidimensional.
\end{prop}

\begin{proof}
If $A \in S_n$, then the claim is trivial, and so we assume $A \notin S_n$.

Suppose $w \in \Perm(A)$.  We will first show that $1 \oplus w \in \Perm(1 \oplus A)$.  From the definition of $1 \oplus A$, we compute \[\rk_{1 \oplus A}(i,j) = 
\begin{cases}
1 & i=1 \mbox{ or } j=1\\
\rk_A(i-1,j-1)+1 & \mbox{otherwise}.
\end{cases}
\]

Thus
\begin{equation}\label{eq:perm_1}
    \text{$w > A$ if and only if $1 \oplus w > {1 \oplus A}$.}
\end{equation}
Suppose there exists a $v$ such that $1 \oplus w \ge v > 1 \oplus A$.  From the equations
\[
    \rk_{1 \oplus w}(1,j) \le \rk_{v}(1,j) \le \rk_{1 \oplus A}(1,j)
\] and \[
 \rk_{1 \oplus w}(1,j) = 1 = \rk_{1 \oplus A}(1,j)
\]
for all $j \in [n+1]$, we see that the first row of $v$ is the same as the first row of $1 \oplus A$, which is the same as the first row of $1 \oplus w$ - specifically, the $1^{st}$ unit row vector.  Similarly, the first column of all three matrices $1 \oplus w$, $v$, and $1 \oplus A$ must be the $1^{st}$ unit column vector.  

Thus $v = 1 \oplus v'$ for some $v' \in S_n$.  By (\ref{eq:perm_1}), $w \geq v'>A$.  By definition of $\Perm(A)$, the inequality $w \geq v' > A$ implies $w=v'$, and so $1 \oplus w = 1 \oplus v'=v$.  Hence, $1 \oplus w \in \Perm(1 \oplus A)$.

In the other direction, fix some $u \in \Perm(1 \oplus A)$.  We claim that $u = 1 \oplus u'$ for some $u' \in S_n$.  It suffices to show that $u(1) = 1$.

Suppose for contradiction that $u(1) = j$ for some $j>1$ and that $i$ is the least index so that $u(i)<j$.  Let $\nu = ut_{1,i}$, where $t_{1,i}$ is the transposition $(1i)$.  That is, in matrix form, $\nu$ is the matrix obtained from $u$ by exchanging rows $1$ and $i$. 

We claim that $1 \oplus A \leq \nu <u$, which will contradict the assumption $u \in \Perm(1 \oplus A)$.  Note that $\rk_\nu(\alpha,\beta) = \rk_u(\alpha,\beta)$ unless $\alpha \in [i-1]$ and $\beta \in [u(i), j-1]$, in which case $\rk_u(\alpha,\beta) = \rk_\nu(\alpha,\beta)-1$ (see \Cref{ex:permBij} for a visualization). Thus $\nu<u$.  Because $i$ was chosen minimally, $\rk_u(\alpha,\beta) = 0$ for all $\alpha \in [i-1]$ and $\beta \in [u(i), j-1]$, and so $\rk_\nu(\alpha,\beta) = 1$ for all such $(\alpha, \beta)$.  But $\rk_{1 \oplus A}(\alpha, \beta) \ge 1$ for all $\alpha,\beta$.  It follows from $\rk_\nu(\alpha,\beta) = \rk_u(\alpha,\beta) \leq \rk_{1 \oplus A}(\alpha, \beta)$ unless $\alpha \in [i-1]$ and $\beta \in [u(i), j-1]$ together with $\rk_\nu(\alpha,\beta) = 1 \leq \rk_{1 \oplus A}(\alpha,\beta)$ for $\alpha \in [i-1]$ and $\beta \in [u(i), j-1]$ that $1 \oplus A \leq \nu$, completing the claim.

Hence our arbitrary $u \in \Perm(1 \oplus A)$ has the form $u = 1 \oplus u'$ for some $u' \in S_n$. 
 It remains to show that $u' \in \Perm(A)$. 
 By (\ref{eq:perm_1}), $u' >A$.  Hence, if $u' \notin \Perm(A)$, then there exists some $\tilde{u} \in S_n$ satisfying $A < \tilde{u}<u'$.  But then, by (\ref{eq:perm_1}), $1 \oplus A < 1 \oplus \tilde{u}<u$, contradicting the assumption $u \in \Perm(1 \oplus A)$.
\end{proof}

\end{document}
